\section{One central-induction certificate framework}\label{sec:unified-induction}
The following subsections are blocks of one restriction-and-trace system.
The Gram data and representation multiplicities are fixed first. Character
labels, reciprocity, Galois orbits, indicators, and modular identities then
add equations or remove impossible columns from that same system.
Each block states its own scope: Frobenius--Perron positivity is used only
in the pseudo-unitary blocks, whereas the signed block retains negative
spherical dimensions. A Farkas vector certifies infeasibility only after
all low frames, character assignments, and permitted indicator choices
have been exhausted. No particular block is mandatory when a shorter
inherited certificate already proves the required exclusion.

\subsection{Central induction with representation multiplicities}\label{sec:weighted}
Assume in this section that $\mathcal C$ is pseudo-unitary, with its
positive spherical structure. Let $F:Z(\mathcal C)\to\mathcal C$ be
forgetful and $I$ its adjoint. Put
\begin{equation}\label{eq:gram}
 B_{ji}=[FI(b_j):b_i]
 =\sum_{a,k}N_{aj}^{k}N_{k,a^*}^{i}.
\end{equation}
This is the induction Gram matrix. In a commutative ring it agrees with
the regular Casimir. In a noncommutative ring one must use
\eqref{eq:gram}, not silently substitute the Casimir.

For a central simple $Z$, write $v_Z=([F(Z):b_i])_i$. Adjunction gives
\begin{equation}\label{eq:fullgram}
 B=\sum_{Z\in\Irr Z(\mathcal C)}v_Zv_Z^{\mathsf T}.
\end{equation}
The summands of $I(\one)$ are indexed by irreducible representations
$E$ of $R\otimes\CC$:
\begin{equation}\label{eq:weightedlow}
 I(\one)=\bigoplus_E r_E A_E,
 \quad (v_{A_E})_0=r_E,
 \quad \dim A_E=D/f_E,
 \quad \theta_{A_E}=1.
\end{equation}
These are the formal-codegree and induction statements of
\cite[Theorems~2.5,2.13]{Ostrik3}.

\begin{theorem}[Multiplicity-weighted low induction]\label{thm:weighted}
For each $E$ let $v_E\in\ZZ_{\ge0}^6$ satisfy
\begin{equation}\label{eq:weightedconstraints}
 (v_E)_0=r_E,\qquad
 d\cdot v_E=D/f_E,\qquad
 \sum_E r_Ev_E=B_{0,*}.
\end{equation}
A necessary condition for a pseudo-unitary categorification is that at
least one such collection has a nonnegative integral Gram completion
\begin{equation}\label{eq:remainder}
 Q=B-LL^{\mathsf T}=\sum_{Z\not\subset I(\one)}v_Zv_Z^{\mathsf T},
 \qquad L=(v_E)_E,
\end{equation}
and satisfies the central-indicator necessary conditions of
\cite{DP4,DP5}. In particular, $Q$ is entrywise nonnegative and positive
semidefinite. The columns of $L$ occur once each in $LL^{\mathsf T}$;
they are not repeated $r_E$ times.

For a noncommutative rank-six ring the low unit multiplicities are
$1,1,2$. Equations~\eqref{eq:weightedconstraints} and
\eqref{eq:remainder}, with these multiplicities, give a finite exhaustive
low-column search.
\end{theorem}
\begin{proof}
The first two identities of \eqref{eq:weightedconstraints} are
\eqref{eq:weightedlow}. Restricting $I(\one)$ proves the third.
Separating the distinct simples occurring in $I(\one)$ from the sum
\eqref{eq:fullgram} proves \eqref{eq:remainder}. This also explains why
$r_E$ weights the row-sum equation but not the number of copies of a
column in $L$.

Finiteness follows from
$(v_E)_i^2\le B_{ii}$ and
$r_E(v_E)_i\le B_{0i}$. Thus each entry belongs to an explicitly bounded
integer interval. All dimension equalities are checked in the exact
number field. The usual trace identities involve only restriction
multiplicities and categorical dimensions, so their proofs and the
finite relative-trace construction do not require commutativity of $R$
once the correct low data and Gram matrix are supplied. Finally a
noncommutative semisimple rank-six algebra with a dimension character
has type \eqref{eq:wedderburn}, giving the stated multiplicities.
\end{proof}

In the implementation the one-dimensional character subspace is
computed before selecting a Frobenius--Perron embedding. The unique
matrix-block eigenvalue $c$ of the Casimir gives $f_E=c/2$, and the
identity
\[
 \det(tI-C_R)=(t-c)^4\det(tI-C_R|_{\rm linear})
\]
is checked. A positive gate from Proposition~\ref{prop:inputs}(iv) is
required before a pseudo-unitary exclusion is promoted to a general
one. The five higher-multiplicity noncommutative Haagerup-type tensors
in the census are excluded in this way. The multiplicity-one Haagerup
ring and the group ring of $S_3$ are retained positive controls.

\subsection{Balanced finite-indicator systems}\label{sec:balanced}
The principal positive-dimension test is a single relative-trace
construction. We strengthen it without introducing a new case-specific
obstruction. Throughout this section suppose all dimensions lie in
$K=\QQ$ or $K=\QQ(\sqrt q)$, with $q>1$ squarefree in the latter case.
The low matrix $L$ is one of the possibilities of
Theorem~\ref{thm:weighted}.

\subsubsection{The inherited finite system}
For a simple $b_i$, the Ng--Schauenburg identity \cite{NS} is
\begin{equation}\label{eq:indicator}
 \sum_Z (v_Z)_i\,\dim(Z)\,\theta_Z^k=D\nu_k(b_i).
\end{equation}
The first indicator vanishes for nonunit simples, and the second is zero
for non-self-dual simples and $\pm1$ for self-dual simples. Higher
indicators satisfy the bounds used in the finite-indicator theorem of
\cite{DP5}. We use exactly that theorem, including its permitted
indicator choices and rational inequalities, rather than assuming
unknown indicators have prescribed values.

For a chosen height $h$, take joint normalized relative traces of
$(\theta,\ldots,\theta^h)$ to $K$. The exponent prime support is
controlled by the categorical Cauchy theorem \cite{BNRW}. Let $P$ be
the prime divisors of the norm of $D$, and let $f$ be the conductor of
$K$ (take $f=1$ for $K=\QQ$). The trace vectors are covered by roots
whose orders divide
\begin{equation}\label{eq:orderbox}
 M_h=\prod_{p\in P}p^{\max\{1,v_p(f)\}+\lfloor\log_p h\rfloor},
\end{equation}
together with the zero vector. As proved in \cite{DP5}, a higher
prime-power layer contributes zero to every one of the first $h$
traces. Thus this is not a bound imposed on the actual twist orders.
Ramanujan and quadratic Gauss sums compute the finite trace vectors
exactly. Extra vectors may be retained as a relaxation; they cannot
produce a false exclusion.

An admissible high column has zero unit entry, a nonnegative integral
tail $v$, and $Q-vv^{\mathsf T}\succeq0$. Entrywise Gram bounds and the
algebraic-integrality condition $D/\dim(v)\in\mathcal O_K$ further
restrict it. We also require $\dim(v)$ to be a $d$-number: the center
is braided, so its simple squared norms are $d$-numbers by
\cite[Theorem~1.8]{Ocode}; in the positive spherical structure these
are the squares of the dimensions, and
\cite[Corollary~2.9]{Ocode} removes the square. In a quadratic field
this is the exact condition that the square of the trace divided by
the norm is integral. Denote the finite list by $\mathcal V$, and the finite
joint trace list by $\mathcal T$. Variables $x_{v,\tau}\ge0$ give the
number of high simples with column $v$ and trace type $\tau$; the
linear relaxation allows nonintegral values. Gram identities,
\eqref{eq:indicator}, and the cited higher-indicator inequalities give
an integer linear system
\begin{equation}\label{eq:lp}
 Ax=b,\qquad x\ge0,
\end{equation}
after adding slack variables and clearing denominators. The code
reconstructs this system from the tensor, the low matrix, and the
indicator choices.

\subsubsection{Duality and Galois balance}
Write $v^*_i=v_{i^*}$. Low simples are stable under duality, so the high
part is stable as well. Duality preserves the ribbon twist:
$\theta_{Z^*}=\theta_Z$, not $\theta_Z^{-1}$. Consequently the
actual integer counts satisfy
\begin{equation}\label{eq:dualbalance}
 x_{v,\tau}=x_{v^*,\tau}
 \qquad(v\in\mathcal V,\ \tau\in\mathcal T).
\end{equation}
If a partner column is absent from the list, its variable is zero.

For the quadratic case let $\sigma(\sqrt q)=-\sqrt q$ and define, for
a positive central dimension $\beta\in K$,
\begin{equation}\label{eq:gammadim}
 \gamma(\beta)=\left|\frac{D}{\sigma(D)}\sigma(\beta)\right|,
\end{equation}
where absolute value is at the fixed real embedding. Let $c_\beta$ be
the number of \emph{distinct} low central simples of dimension $\beta$.
The low multiplicity $r_E$ in $I(\one)$ is not included in this count.

\begin{lemma}[Dimension-count balance]\label{lem:balance}
Galois symmetry permutes the central simples with dimension
$\beta$ onto those with dimension $\gamma(\beta)$. Hence their high
counts satisfy
\begin{equation}\label{eq:galbalance}
 \sum_{\dim(v)=\beta,\tau}x_{v,\tau}
 -\sum_{\dim(v)=\gamma(\beta),\tau}x_{v,\tau}
 =c_{\gamma(\beta)}-c_\beta.
\end{equation}
Moreover $\gamma^2(\beta)=\beta$.
\end{lemma}
\begin{proof}
Extend $\sigma$ to a cyclotomic field of the center. Its normalized
$S$-matrix is $s=S/D$, since the center has global dimension $D^2$.
Signed Galois symmetry gives
\[
 \sigma(s_{0Z})=\epsilon_Zs_{0,\widehat\sigma(Z)},
 \qquad\epsilon_Z\in\{1,-1\}.
\]
The entry $s_{0Z}$ equals $\dim(Z)/D$. All central dimensions are
positive. Taking absolute values therefore gives
$\dim\widehat\sigma(Z)=\gamma(\dim Z)$, proving the permutation
statement. Subtract the known low counts to obtain
\eqref{eq:galbalance}. Finally the map
$u\mapsto (D/\sigma D)\sigma(u)$ squares to the identity. Taking
absolute values commutes with changing a real sign, so the positive
version squares to the identity as well. See also the squared-dimension
form of Galois symmetry in \cite[Equation~(2.4)]{MY}.
\end{proof}

The factor $D/\sigma(D)$ and the absolute value in
\eqref{eq:gammadim} are both essential. Omitting either does not give a
valid positive dimension permutation in general.

\begin{theorem}[Balanced finite-indicator criterion]\label{thm:balanced}
Append \eqref{eq:dualbalance} and, when quadratic,
\eqref{eq:galbalance} to the finite system of \cite{DP5}.
If for every admissible low matrix and every allowed second-indicator
choice the resulting system has an exact Farkas certificate
\begin{equation}\label{eq:farkas}
 A^{\mathsf T}y\ge0,\qquad b^{\mathsf T}y<0,
\end{equation}
then there is no pseudo-unitary categorification. Coupled with a valid
general-to-positive reduction, this excludes general complex
categorification.

At a fixed height and with the same admissible low and high columns,
this criterion contains the former unbalanced finite-indicator test:
every old separating functional extends by zero on the appended rows.
\end{theorem}
\begin{proof}
An actual category supplies integral nonnegative counts satisfying the
inherited system and both new balances, by duality and
Lemma~\ref{lem:balance}. But \eqref{eq:farkas} would imply
$0>b^{\mathsf T}y=x^{\mathsf T}A^{\mathsf T}y\ge0$.
The last assertion follows by appending zero coordinates to $y$.
\end{proof}

This containment justifies using the balanced theorem as the common
finite-indicator framework. It does not imply that every earlier
non-indicator criterion is redundant. In particular, the cheap
arithmetic filters and signed reductions remain useful.

\subsubsection{An integral refinement and its scope}
For a high-column list $v_1,\ldots,v_s$, enumerate nonnegative integers
$m_j$ with
\begin{equation}\label{eq:integergram}
 Q=\sum_{j=1}^s m_jv_jv_j^{\mathsf T}.
\end{equation}
This is finite because each nonzero $v_j$ has a positive diagonal entry.
For each frame append
$\sum_\tau x_{v_j,\tau}=m_j$ to the balanced system. Exhaustive
infeasibility for every frame would be a further valid obstruction.
We implement this refinement principally to diagnose the frontier.
The feasible point in Section~\ref{sec:frontier} survives it. We do
not count this unproductive refinement as an independent successful
exclusion criterion.

\subsection{Cyclic-character compatibility for tensor powers}\label{sec:rotation}
The preceding indicator system bounds different higher indicators
separately. There is additional compatibility: several of them are
traces of powers of a single finite-order linear operator. This gives
extra linear equations inside the same exact-certificate framework.

\begin{lemma}[Block rotation]\label{lem:blockrotation}
In a pivotal fusion category, let $E_X^{(n)}$ be cyclic rotation on
$\operatorname{Hom}(1,X^{\otimes n})$. For every positive divisor $r$ of $n$,
\begin{equation}\label{eq:blockrotation}
 (E_X^{(n)})^n=\mathrm{id},\qquad
 \Tr((E_X^{(n)})^r)=\nu_{n/r}(X^{\otimes r})
   =\sum_i[X^{\otimes r}:X_i]\nu_{n/r}(X_i).
\end{equation}
Here the trace is the ordinary trace on the displayed Hom space.
\end{lemma}
\begin{proof}
Pivotal strictification preserves the rotation operators. In the strict
pivotal setting, equation (5.4) in \cite{NSpiv} states that consecutive
rotations of adjacent tensor blocks equal the rotation of their tensor
product. Induction gives
$(E_X^{(n)})^r=E_{X^{\otimes r},X^{\otimes(n-r)}}$, exactly the identity
used in the proof of \cite[Theorem~5.1]{NSpiv}. If $r\mid n$, group the
factors into $n/r$ consecutive blocks of length $r$. This identifies the
operator with one rotation on the blocks. Taking traces gives the first
indicator equality; the final equality is additivity of higher
indicators on a direct sum. For $r=n$ it is the identity operator, with
trace $[X^{\otimes n}:1]$.
\end{proof}

Fix the dimension field $K$, the low columns, and the second indicators
as in Section~\ref{sec:balanced}. For a positive integer $n$, let
$\mathcal U_{K,n}$ be the finite set of vectors
\[
 (P_K(\zeta^r))_{r\mid n,\ r<n},\qquad \zeta^n=1,
\]
where $P_K$ is normalized relative trace from a common cyclotomic
extension to $K$. This definition is independent of enlarging the common
extension. The vectors are computed by the same Ramanujan and quadratic
Gauss sums used for the center twists, but their orders now divide $n$
exactly. In particular no unbounded-order approximation occurs in this
auxiliary set.

\begin{theorem}[Joint tensor-rotation criterion]\label{thm:rotation}
In a pseudo-unitary categorification with simple dimensions in $K$, for
each simple $X$ and each $3\le n\le h$ there are nonnegative integers
$y_{X,n,u}$, $u\in\mathcal U_{K,n}$, such that
\begin{align}
 \sum_u y_{X,n,u}&=[X^{\otimes n}:1],\label{eq:rotation-dim}\\
 \sum_u y_{X,n,u}u_r
 &=\sum_i[X^{\otimes r}:X_i]P_K(\nu_{n/r}(X_i))
       \qquad(r\mid n,\ r<n).\label{eq:rotation-characters}
\end{align}
Together with the center trace identities, these are linear equations
in the existing high-center counts and the new $y$ variables after
multiplication by $D$ and expansion in a rational basis of $K$.
Infeasibility of the nonnegative real relaxation is an obstruction to
pseudo-unitary categorification. Combined with any proved
sphericalization or pseudo-unitarity reduction, it is a general
categorification obstruction.
\end{theorem}
\begin{proof}
The operator in Lemma~\ref{lem:blockrotation} is diagonalizable over
$\CC$ and has $n$th roots of unity as eigenvalues. Group its eigenvalues,
with multiplicity, by their joint relative-trace vectors. This gives
nonnegative integer counts satisfying \eqref{eq:rotation-dim} and
\eqref{eq:rotation-characters}. Use \eqref{eq:indicator} to express each
indicator on the right. Its denominator is the known scalar $D$, while
its numerator is the known low contribution plus a linear combination
of the high-center variables. Since $P_K$ is $K$-linear, multiplying by
$D$ gives linear equations with coefficients in $K$. Clearing rational
denominators gives an integer system. An integer solution would also be
a nonnegative real solution; an exact separating functional therefore
proves the claimed obstruction.
\end{proof}

For clarity, let $\ell_i$ be the dimension of the known low part of
$I(X_i)$, with $\ell_0=D$, and let $(v,\tau)$ index a high variable.
The equation appended for a divisor $r$ is exactly
\begin{equation}\label{eq:rotation-row}
 D\sum_u y_{X,n,u}u_r
 -\sum_{v,\tau}x_{v,\tau}\Bigl(\sum_{i>0}[X^{\otimes r}:X_i]v_i\Bigr)
                 \dim(v)\tau_{n/r}
 =\sum_i[X^{\otimes r}:X_i]\ell_i.
\end{equation}
This formula specifies the implemented rows without a graphical
normalization ambiguity. The unit indicator is one, and unit-free high
columns have unit entry zero.

The criterion strengthens the existing system by adding equations and
variables with nonnegative constraints; no previously valid inequality
is dropped. Old dual certificates extend by zero. Its height-six
application excludes $\RR{2}{55}$, the self-dual Klein-four quadratic
ring with two noninvertibles of dimension $2+\sqrt6$. The complete ten
low-column alternatives and every permitted second-indicator choice
are checked. The two-Frobenius--Perron-dimension reduction supplies the
general scope. No separate theorem declaring this particular ring
impossible is used by the program.

The cyclic-four ring $\RR{2}{4}$ has the same dimension type but does not
satisfy that exclusion certificate. Section~\ref{sec:frontier} records
an exact integer-count feasible point for the stronger system. Thus
neither dimension-type identification nor a failed numerical optimization
can replace full-tensor testing.

\subsection{Character-aware central induction}\label{sec:character-induction}
The dimension-only induction matrix forgets an important part of the
formal-codegree correspondence: each of its low columns has a specified
\emph{character}, not only a specified dimension. Retaining these
characters adds exact modular reciprocity and a joint Galois-spectrum
condition to the same finite-indicator system. No classification of
modular data, numerical eigenvalue decision, or enumeration of center
fusion rings is required.

\subsubsection{The character attached to a low column}
Let $\mathcal C$ be pseudo-unitary, let $R=K_0(\mathcal C)$ be
commutative of rank $n$, and write $d_i=\FPdim(X_i)$ and
$D=\sum_i d_i^2$. For a character $\chi:R\to\CC$, put
\[
 f_\chi=\sum_i\chi(X_i)\chi(X_i^*),\qquad
 \beta_\chi=D/f_\chi.
\]
Use the unnormalized modular matrix $S$ of $Z(\mathcal C)$, so that
$S_{1W}=\dim W$. The correspondence preceding
\cite[Theorem~2.13]{Ostrik3} identifies the simple summand $Z_\chi$ of
$I(1)$ by the equality of characters
\begin{equation}\label{eq:central-character}
 \chi(FW)=\frac{S_{Z_\chi,W}}{\beta_\chi}
 \qquad(W\in\Irr Z(\mathcal C)).
\end{equation}
Indeed, pull back $\chi$ along $F$. The resulting character of the
modular fusion ring is its unique normalized $S$-row, and
\cite[Theorem~2.13]{Ostrik3} identifies the row label, its dimension
$\beta_\chi$, and its multiplicity one in $I(1)$. Thus this statement
uses the character-level correspondence, not an arbitrary choice among
central simples of equal dimension.

For an integral restriction vector $v=(v_i)$ write
$d(v)=\sum_i v_i d_i$ and $\chi(v)=\sum_i v_i\chi(X_i)$.
If $a_\chi$ is the restriction vector of $Z_\chi$, symmetry of $S$
gives
\begin{equation}\label{eq:character-reciprocity}
 \boxed{\ \beta_\chi\chi(a_\psi)
             =\beta_\psi\psi(a_\chi)\ }.
\end{equation}
The unit column is labelled by the Frobenius--Perron character. When
formal codegrees repeat, every dimension-compatible assignment of the
remaining distinct characters must be considered. Selecting one
assignment without exhausting the others would not be an obstruction.

\subsubsection{Joint Galois closure}
The next elementary spectral observation is useful well beyond the
present induction problem.
\begin{lemma}[Joint spectral closure]\label{lem:joint-spectrum}
Let $B$ be a commutative fusion ring and $b_1,\ldots,b_s\in B$.
The multiset
\[
 \{(\varphi(b_1),\ldots,\varphi(b_s)):
                    \varphi\in\Irr(B\otimes\overline\QQ)\}
\]
is invariant under every automorphism of $\overline\QQ$ over $\QQ$.
Consequently a finite set containing all these tuples must contain
every Galois conjugate of each tuple that actually occurs.
\end{lemma}
\begin{proof}
The multiplication equations of $B$ have integral coefficients. Composing
a character with a field automorphism gives another character, and the
inverse automorphism makes this operation a permutation. Equivalently,
the commuting integer multiplication matrices have a Galois-invariant
joint spectrum. Restriction to the specified elements preserves the
multiset, including repeated restricted tuples.
\end{proof}

For a nonzero prospective restriction vector $v$ define
\begin{equation}\label{eq:central-eigen-tuple}
 \Lambda(v)=
 \left(\frac{\beta_\chi\chi(v)}{d(v)}\right)_{\chi\in\Irr R}.
\end{equation}
If $v=FW$ for a central simple $W$, then
\eqref{eq:central-character} and symmetry give
\[
 \frac{\beta_\chi\chi(v)}{d(v)}
       =\frac{S_{W,Z_\chi}}{\dim W}.
\]
These are the joint eigenvalues of the integer fusion matrices of
$Z_\chi$ evaluated at the character labelled by $W$. In particular,
Lemma~\ref{lem:joint-spectrum} applies to the \emph{whole tuple}.
Checking its coordinates separately would be weaker.

Fix a candidate low matrix $L$. Let $\mathcal H_L$ be any finite set
containing all possible high restriction columns for that low matrix;
the Gram and algebraic-integrality construction in
Section~\ref{sec:inputs} and the finite-indicator implementation supplies
such a set. Let $\mathcal A_L$ be its union with the low columns, and set
\[
 \mathcal T_L=\{\Lambda(v):v\in\mathcal A_L\}.
\]
Choose a finite Galois field $K$ containing all character values and all
$d_i$. Every coordinate of every tuple belongs to $K$. A column $v$ is
\emph{forbidden} if $\sigma(\Lambda(v))\notin\mathcal T_L$ for some
$\sigma\in\operatorname{Gal}(K/\QQ)$. A forbidden low column rejects
$L$; a forbidden high column must have multiplicity zero. Notice that
\eqref{eq:central-eigen-tuple} uses the fixed character order and
$D/f_\chi$, so this test is independent of the permutation used to
label equal-dimensional low columns.

\begin{theorem}[Character-aware finite-indicator obstruction]
\label{thm:character-induction}
For every low matrix of a pseudo-unitary categorification with
commutative Grothendieck ring:
\begin{enumerate}[label=(\roman*)]
\item there is a dimension-compatible bijection from its columns to
$\Irr R$ satisfying \eqref{eq:character-reciprocity};
\item none of its low columns is forbidden by the joint Galois test;
\item the finite-indicator system remains feasible after all variables
supported on forbidden high columns are set to zero.
\end{enumerate}
Therefore the absence of an assignment in (i), failure of (ii), or an
exact Farkas obstruction to (iii), for every candidate low matrix and
every permitted second-indicator choice, excludes a pseudo-unitary
categorification.
\end{theorem}
\begin{proof}
The formal-codegree correspondence proves (i). Every actual central
restriction belongs to $\mathcal A_L$. Lemma~\ref{lem:joint-spectrum},
applied to the low central simples, proves (ii) and the asserted zero
multiplicities. The finite-indicator variables group actual central
simples by restriction and joint relative twist trace. Every variable
with a forbidden restriction is therefore zero; the remaining actual
multiplicities satisfy the inherited equations and inequalities.
After slack variables are inserted these have the form $Ax=b$, $x\ge0$.
A rational vector $y$ with $A^{\mathsf T}y\ge0$ and $b^{\mathsf T}y<0$
is impossible, because $b^{\mathsf T}y=x^{\mathsf T}A^{\mathsf T}y\ge0$.
All ordinary low-matrix rejections remain applicable before these
additional tests.
\end{proof}

This is a strengthening of the existing finite-indicator framework,
not a competing list of special forbidden rings. The implementation
uses exact splitting fields of degree at most sixteen generated by
quadratic eigenvalues of one or two ring generators. This is an
implementation eligibility bound, not a hypothesis of the theorem.
Every proposed character is checked against all multiplication
equations; the complete set has cardinal $n$. Field automorphisms are
constructed by splitting the primitive-element polynomial over $K$ and
checking that all of its distinct roots have been obtained. The verifier
reconstructs all labellings, all possible high columns, the forbidden
set, and every rational separating vector. It never calls the numerical
optimizer used to discover a vector.

\subsubsection{General categorification and the multiplicity-two completion}
Two illustrative tensors give a compact reproducibility summary:
\begin{center}
\begin{tabular}{lrrl}\toprule
Ring & Low matrices & Reciprocity survivors & General reduction\\\midrule
$\RR{2}{20}$ & 77 & 2 & two FP dimensions\\
$\RR{2}{41}$ & 84 & 2 & two exhaustive signed tables\\\bottomrule
\end{tabular}
\end{center}
For $\RR{2}{20}$, both surviving matrices and both permitted second
indicators have exact height-six obstructions. The two-FP-dimension
Galois reduction of Proposition~\ref{prop:inputs} makes this a general
exclusion. For $\RR{2}{41}$, one reciprocity survivor is already rejected
by the inherited height-six equations. For the other, there are $104$
possible high columns; joint Galois closure forbids $66$, and the
restricted height-six system has an exact separating vector.

For this latter tensor the complete oriented signed enumeration has
exactly two tables. Both are sign gauges of the original tensor, with
sign vectors $(1,-1,-1,1,1,1)$ and $(1,1,1,1,1,1)$ in census order.
The ordinary spherical-character test excludes every character outside
the FP Galois orbit by a zero dimension or unequal dimensions on a dual
pair. That orbit is excluded by Theorem~\ref{thm:character-induction}.
The signed-sphericalization reduction already established in the
preceding sections therefore excludes \emph{general} categorification,
without assuming the conjectural existence of pivotal structures.
These statements are recomputed from the tensors in the certificate,
including the completeness of the oriented enumeration.

Table~\ref{tab:character-applications} records $\CharacterExampleCount$ certified examples
of this common character-aware mechanism. Their stored proof trees
include the complete sphericalization or Galois reduction required for
general nonexistence. In particular, the signed calculation for
$\RR{3}{132}$ is detailed in Section~\ref{sec:modular-frontier}.
The certificate controls recompute these applications without a
numerical optimizer. This illustrative collection is not a separate
allocation of the entire census: the current sequential counts,
eligibility outcomes and bounded-search diagnostics appear in
Section~\ref{sec:verification} and its machine-readable ledger.

\begin{table}[ht]\centering
\begin{tabular}{cr}\toprule
Exact multiplicity & Example certificates\\\midrule
2 & 2\\
3 & 2\\
4 & 6\\
5 & 3\\
6 & 10\\
\midrule Total & 23\\\bottomrule
\end{tabular}
\caption{Illustrative character-aware induction certificates for Theorem~\ref{thm:character-induction}, generated from the supplied exact proof examples. These are examples of the mechanism, not a disjoint allocation of the current census.}\label{tab:character-applications}
\end{table}



The multiplicity-three applications are detailed in
Section~\ref{sec:modular-frontier}. The complete consecutive classification
range now reaches multiplicity three by the stabilizer-rotation argument
of Theorem~\ref{thm:frontier-noncategorifiable}. The counting bound eight is
not presented as a categorification bound.

\subsubsection{Modular equations retained inside the indicator system}
\label{sec:modular-frontier}
We describe one further strengthening of
Theorem~\ref{thm:character-induction}. It uses the same high-column and
relative-trace variables; it does not require enumerating an unknown
center fusion ring. Its purpose here is also to specify exactly which
stronger necessary conditions the former multiplicity-three residual
satisfies before the stabilizer constraint is imposed.

Assume that all ring characters and positive dimensions lie in a real
field $K$ of degree at most two, and let $P$ be the normalized relative
trace from a common cyclotomic field to $K$. For a central simple $W$
write $v_W=FW$, $d_W=d(v_W)$, and
\[
 H_{aW}=S_{Z_a,W}=\beta_a\chi_a(v_W).
\]
Here $Z_a$ ranges over the labelled low objects. Thus $H_{aW}$ denotes an entry of the unnormalized matrix $S$.

\begin{proposition}[Joint dimension, twist and low modular constraints]
\label{prop:modular-frontier}
In addition to the preceding induction conditions, every pseudo-unitary
categorification satisfies the following statements.

Every coordinate of $\Lambda(v_W)$ is an algebraic integer. Under a
Galois automorphism $\sigma$, the pair consisting of a simple dimension
and its low fusion-eigenvalue tuple is permuted according to
\begin{equation}\label{eq:augmented-galois}
 (d_W,\Lambda(v_W))\longmapsto
 \left(\left|\frac{D}{\sigma(D)}\sigma(d_W)\right|,
                         \sigma(\Lambda(v_W))\right).
\end{equation}
This permutation preserves every joint relative twist-trace type
\[
 (P(\theta_W),P(\theta_W^2),\ldots,P(\theta_W^h)).
\]
Finally, for any two low labels $a,b$,
\begin{equation}\label{eq:low-modular-pair}
 \sum_W P(\theta_W)H_{aW}H_{bW}=D S_{Z_a,Z_b}.
\end{equation}
Thus the finite indicator system may impose equality of multiplicities
within each joint orbit in \eqref{eq:augmented-galois}, separately for
each trace type, and may impose all equations
\eqref{eq:low-modular-pair}. A missing orbit member forbids a candidate
column; an exact infeasibility certificate for the enlarged system is a
pseudo-unitary obstruction.
\end{proposition}
\begin{proof}
The entries of $\Lambda(v_W)$ are eigenvalues of integer fusion
matrices, so they are algebraic integers. Use the canonical normalized
modular representation of the center, namely $S/D$ and $T$. Its Galois
symmetry gives
\[
 \sigma(S_{0W}/D)=\epsilon_\sigma(W)S_{0,\widehat\sigma W}/D,
 \qquad
 \Lambda(v_{\widehat\sigma W})=\sigma(\Lambda(v_W)),
\]
where $\epsilon_\sigma(W)\in\{1,-1\}$; see
\cite[equations~(2.12)--(2.13)]{BNRW}. Positivity of all simple
Frobenius--Perron dimensions gives exactly
\eqref{eq:augmented-galois}, including its absolute value.

For this canonical center normalization, the Galois symmetry of $T$
is $\theta_{\widehat\sigma W}=\sigma^2(\theta_W)$, by
\cite[Theorem~2.9(b)]{BNRW}. A square of a Galois automorphism fixes
$K$, since its degree is at most two. The normalized trace $P$ is
invariant under automorphisms over $K$, and cyclotomic Galois groups
are abelian. Consequently
$P(\theta_{\widehat\sigma W}^{\,j})=P(\theta_W^j)$ for every $j$.
The simultaneous dimension-and-eigenvalue permutation therefore gives
multiplicity balances within each trace type, including the low objects
in the twist-one type. This statement concerns multiplicities of tuples;
it does not assert that a restriction vector itself is Galois invariant.

The positive Gauss sum of the center is $D$. Its modular relation is
\[
 D S_{Z_a,Z_b}
    =\theta_{Z_a}\theta_{Z_b}\sum_W\theta_W S_{Z_a,W}S_{W,Z_b},
\]
by \cite[equation~(2.6)]{BNRW}. Both low twists are one. All the other
coefficients in this equation belong to $K$, by the character formula.
Applying $P$ proves \eqref{eq:low-modular-pair}. The actual central
simples satisfy the resulting rational equations and nonnegative
multiplicity constraints. An exact Farkas obstruction therefore applies
as in Theorem~\ref{thm:character-induction}.
\end{proof}

The code reconstructs all finite Galois conjugates, checks algebraic
integrality by integer characteristic polynomials, and adds coefficient
rows over $\QQ$. This is one refinement of the common induction system,
not an independently allocated family of forbidden rings. In particular,
the final general exclusion of $\RR{3}{132}$ already has a certificate
using Theorem~\ref{thm:character-induction}; no stronger criterion is
credited when an inherited one suffices.

\subsubsection{Exact multiplicity-three outcome and the higher scan}
The general certificate for $\RR{3}{132}$ has the following finite
structure. Its positive type is $(1,1,d,d+1,d+1,d)$, with
$d=4+\sqrt{21}$, and its global dimension is $168+36\sqrt{21}$.
There are $172$ low matrices. Modular reciprocity rejects $168$;
the other four have exact height-six indicator certificates after the
joint-spectrum exclusions. The general sphericalization reduction
exhausts sixteen oriented signed tables. Eight are sign gauges of the
original table; for the other eight, all character branches are rejected
by the checked norm or zero-dimension tests. These are certificate
statistics reconstructed from the full tensor. They are not premises
selected by its identifier.

Corollary~\ref{cor:inversion-five} supplies the unitary model for
$\RR{3}{218}$. Thus the only remaining exact multiplicity-three tensor
is $\RR{3}{215}$. The strengthened rational system above has an exact
feasible point for that tensor: $230$ equations in $199$ nonnegative
integer variables, including all indicator-slack coordinates, and formal central
rank $22$. Its low frame, character labelling, second indicators and
all rational coordinates are stored in
\filename{systems/integral_frontier_215.json}. The verifier reconstructs
the equations and checks the equalities exactly. That finite point alone
does not specify individual twists or a complete associative center
fusion tensor. Section~\ref{sec:cyclic-eight-center} now supplies a
compatible exact completion of those data. Neither the finite point nor
that completion proves categorifiability or noncategorifiability.

The common framework is applied at every multiplicity in the extended
census. Structure-specific pseudo-unitary certificates remain separate
from general exclusions. In particular the non-positive branches of
$\RR{4}{375}$ and $\RR{4}{446}$ are closed by
Section~\ref{sec:dimension-order}; a positive-dimension obstruction alone
would not suffice. The surviving pseudo-unitary-only exclusions are
$\RR{7}{3582}$ and $\RR{7}{3812}$.
The former structure-only case $\RR{4}{162}$ is now generally excluded
by Theorem~\ref{thm:rank3-subring}.
The complete general ledger and its distinct structure fields, rather
than historical bounded-search outcomes, state the classification.

\subsection{The dimension order and signed central traces}\label{sec:dimension-order}
The positive-dimension induction system has two extensions that are useful
when an exhaustive sphericalization search leaves a non-positive
spherical character. The first sharpens an arithmetic condition without
adding any central variables. The second retains the algebraic part of
the same trace system, but discards every inequality that depends on
positivity of categorical dimensions.

\subsubsection{The actual order of simple dimensions}
Let $d=(1,d_1,\ldots,d_{r-1})$ be a dimension character under
consideration, and let $K=\QQ(d_i)$ and $\delta=\sum_i d_i d_{i^*}$.
The finitely generated group
\[
 \mathcal O_d=\sum_i\ZZ d_i\subset K
\]
is an order: it contains $1$, consists of algebraic integers, and is
closed under multiplication because the $d_i$ satisfy the multiplication
equations with integer coefficients. It need not be the maximal order
of $K$.

\begin{proposition}[Dimension-order obstruction]\label{prop:dimension-order}
If a fusion ring $R$ has a spherical categorification with simple
dimensions $d_i$, then for every formal codegree $f$ of $R$,
\begin{equation}\label{eq:dimension-order}
 f\in K,\qquad \delta/f\in\mathcal O_d.
\end{equation}
The same conclusion holds for the dimensions of selected lifts in the
canonical sphericalization of an arbitrary categorification of $R$.
Thus the order-membership test may be applied to each dimension-character
branch of a complete signed-sphericalization search.
\end{proposition}
\begin{proof}
In the spherical case the central simple $A_E$ attached to an irreducible
representation of $R$ with codegree $f$ has dimension $\delta/f$, by
\cite[Theorem~2.13]{Ostrik3}. Its forgetful image is an object of the
original category. Its dimension is therefore a nonnegative integral
linear combination of the $d_i$, in particular an element of
$\mathcal O_d$. Since $\delta$ is nonzero, this also puts $f$ in $K$.

For an arbitrary categorification use its sphericalization, of dimension
$2\delta$. An irreducible representation of $R$ lifts through the
forgetful quotient to one with formal codegree $2f$; both selected and
opposite basis lifts contribute to that codegree. The same central
induction theorem gives dimension $(2\delta)/(2f)=\delta/f$.
All simple dimensions of the sphericalization are among $\{d_i,-d_i\}$,
so forgetting this central object gives an integral linear combination
of the selected dimensions. This proves the assertion without assuming
that the selected lifts themselves form a fusion subcategory.
\end{proof}

The test is finite and exact. Express the $d_i$ in a rational power basis
of $K$, clear one common denominator, and put their coordinate lattice
in Hermite normal form. Membership of $\delta/f$ is then an integer
linear-algebra question. A negative certificate can equivalently give a
rational linear functional $\ell:K\to\QQ$ satisfying
\[
 \ell(d_i)\in\ZZ\quad\hbox{for all }i,
 \qquad \ell(\delta/f)\notin\ZZ.
\]
The checker verifies the defining polynomial of $f$, reconstructs the
lattice and the functional, and checks both displayed conditions. A
failure on one character component is not a noncategorification proof
unless all possible sphericalization and character branches are closed.
In particular a failed character of a positively realized ring cannot
be promoted to an exclusion of that ring.

For illustration, the complete certificate for $\RR{4}{375}$ has four
oriented signed tables, all gauges of the original tensor. Its nonzero
spherical components have order $\ZZ[\sqrt5]$ and dimensions
\[
 \delta=70+30\sqrt5\quad\hbox{or}\quad\delta=30-10\sqrt5.
\]
The codegree $4$ gives quotients outside that order in both cases.
The supplied functional certificates, not this identifier or this
displayed calculation, are the inputs to the automated verification.

\subsubsection{The algebraic trace system without positive dimensions}
Suppose a spherical character takes values in a real field $K$ of degree
at most two, has nonzero simple dimensions, and satisfies
$d_i=d_{i^*}$. Write $\delta=\sum_i d_i^2$. The restriction vectors of
central simples remain nonnegative integral vectors, even when some
categorical dimensions are negative. Consequently the Gram identity
\eqref{eq:fullgram}, its finite entry bounds, and the formal-codegree
low-column equalities remain valid. They do not use an inequality
between categorical dimensions.

For a commutative input let $\chi_a$ be a character label of a low simple
$Z_a$ and put $\beta_a=\delta/f_a$. If $v_W$ is a candidate restriction
vector, define
\[
 d_W=d\cdot v_W,\qquad
 H_{aW}=\beta_a\chi_a(v_W),\qquad
 \Lambda_W=(H_{aW}/d_W)_a.
\]
Assume that the full ring character table is defined over $K$. The
character correspondence of Theorem~\ref{thm:character-induction} gives
$H_{aW}=S_{Z_a,W}$, with $S$ unnormalized. The following formulation
uses no absolute value of $d_W$.

\begin{theorem}[Signed central-induction obstruction]\label{thm:signed-traces}
A spherical categorification with the specified character supplies a
nonnegative integral solution of the following finite necessary system.
The unknowns count high central restriction vectors and their joint
relative twist-trace types. They satisfy the Gram equations, the signed
identities
\begin{equation}\label{eq:signed-indicators}
 \sum_W(v_W)_i d_W\theta_W^j=\delta\nu_j(b_i),
\end{equation}
and the cyclic-operator bounds on the right-hand indicators. Low
columns have twist one and dimension $\beta_a$. They admit at least
one complete character labelling satisfying modular reciprocity.

All coordinates of $\Lambda_W$ are algebraic integers. Their joint
Galois types have balanced multiplicities, separately for each relative
twist-trace type. An absent Galois type forbids the corresponding
restriction vector. Finally, if $P$ is the normalized trace from a
common cyclotomic field to $K$, then
\begin{equation}\label{eq:signed-low-modular}
 \sum_W P(\theta_W)H_{aW}H_{bW}=\delta S_{Z_a,Z_b}
\end{equation}
for every pair of low labels.

Exhaustive failure of the system for all bounded low columns, all
compatible character labellings, and all permitted second indicators
excludes the specified spherical character. Exact rational Farkas
certificates suffice for failure of the nonnegative rational relaxation.
\end{theorem}
\begin{proof}
Adjunction proves the Gram identity with nonnegative restriction
multiplicities. Central induction supplies the low dimensions and their
unit multiplicities exactly as in \eqref{eq:weightedlow}; neither
statement assumes pseudo-unitarity. The trace formula
\cite{NS} proves \eqref{eq:signed-indicators} with the actual, possibly
negative, dimensions. Bounds obtained as traces of finite-order cyclic
operators on ordinary complex vector spaces do not assume positivity
of categorical dimensions. Only these algebraic bounds, and not the
positive low-mass inequalities, are retained here.

The finite relative-trace list is justified by the exponent-prime
support theorem \cite{BNRW} and the same relative-trace calculation
used in \eqref{eq:orderbox}, with $\delta$ in place of $D$. The high
columns are bounded by the positive integral Gram matrix, not by the
sign of $d_W$. Their additional algebraic-integrality and $d$-number
conditions follow from the spherical modular center as in the proof of
Theorem~\ref{thm:balanced}. Columns of dimension zero are impossible.

Each $\Lambda_W$ is the restriction of the simultaneous character of
integer fusion matrices of the center to the labelled low objects.
It therefore consists of algebraic integers, and its full Galois orbit
occurs with the corresponding multiplicities. For the normalized
center representation $S/\delta,T$, the permutation of columns sends
$\theta_W$ to $\sigma^2(\theta_W)$; see
\cite[Theorem~2.9]{BNRW}. A square fixes the real field $K$ of degree
at most two, so normalized relative traces of every twist power are
unchanged. This proves the asserted balances of joint eigenvalue types.
It does not identify $d_W$ with its absolute value, nor impose the
positive dimension-orbit rule \eqref{eq:augmented-galois}.

The center has Gauss sum $\delta$. Its modular relation, together with
the twist-one property of low objects, is
$\sum_W\theta_W S_{Z_a,W}S_{W,Z_b}=\delta S_{Z_a,Z_b}$.
All coefficients other than the twists lie in $K$. Taking $P$ proves
\eqref{eq:signed-low-modular}. The actual central simples give integer
counts satisfying all the displayed equations. Expanding them in a
rational basis of $K$ and relaxing to nonnegative rational counts can
only enlarge the solution set. Farkas' lemma therefore proves the
last assertion.
\end{proof}

\paragraph{Exact scope of the implementation.}
The routine first checks the complete signed character table, its
multiplication equations, duality, nonzero dimensions and field degree.
It enumerates low columns using Gram bounds and exact dimension
equalities, without calling a positive-dimension low-mass test. It then
uses signed versions of the inherited indicator rows and adds joint
spectrum and low modular rows. In particular it does not reuse the
absolute-value dimension balances of
Proposition~\ref{prop:modular-frontier}. Only the second-indicator rows depend on the chosen signs, and each such row depends on one sign. The checker reconstructs the all-minus and all-plus endpoint systems and selects the exact normalized row for each sign choice. Independent mixed-sign reconstructions check this optimization; every separating-vector inequality is still verified. The second indicators of
self-dual objects are both allowed unless an independently checked
condition forces their sign.

For $\RR{4}{446}$ the remaining non-positive component is
\[
 d=(1,1,-1,-1,2-\sqrt6,2+\sqrt6),\qquad\delta=24.
\]
There are $290$ low matrices. Reciprocity leaves two, each with two
character labellings; all $32$ choices of second indicators are retained
for each labelling. The $128$ exact separation certificates exclude
these systems. The positive character component has the preceding
character-induction obstruction. The full oriented search has sixteen
gauge tables, so the assembled certificate excludes general
categorification, not just this displayed spherical structure.

\subsubsection{Application and structure-specific conclusions}
Component-level rejection and general nonexistence are different
conclusions. For example, rejecting one character of $\RR{4}{174}$ does
not exclude that ring. Complete proof trees for the order and signed
framework include
\[
 \RR{4}{375},\quad\RR{4}{446},\quad
 \RR{5}{1507},\quad\RR{5}{1508},\quad\RR{5}{1582},\quad
 \RR{6}{2767},\quad\RR{7}{1907}.
\]
For the last five, shorter norm and moment certificates close the
exhaustive signed alternatives, and these are retained rather than
credited to a stronger criterion unnecessarily. The finite applications
are specified by full tensors and exact proof trees.

The separate certificates for $\RR{7}{3582}$ and
$\RR{7}{3812}$ exclude pseudo-unitary and therefore unitary
categorification only. Their general, pivotal and spherical existence
fields remain unknown. This distinction is enforced by the ledger
checker, not left to a reader's interpretation of a single status label.

\subsection{Sphericality and Galois reduction at multiplicity three}
\label{sec:frontier-spherical-reduction}
\begin{proposition}\label{prop:frontier-galois-pu}
Every complex fusion category with based ring $\RR{3}{215}$ admits a
spherical structure and has a pseudo-unitary Galois conjugate.
Consequently general categorifiability of this ring is equivalent to
pseudo-unitary categorifiability. This reduction alone asserts neither
existence nor nonexistence; the final obstruction is proved in
Theorem~\ref{thm:frontier-noncategorifiable}.
\end{proposition}
\begin{proof}
Apply the exhaustive oriented signed-sphericalization construction of
Section~\ref{sec:signed}, using every coefficient in its full parity
interval and imposing no positive-indicator or modulo-four normalization.
There are $1480$ associative signed quotients. The $32$ unit-preserving
sign gauges partition them into $63$ orbits; one orbit consists of the
$32$ gauges of the original fusion tensor. On representatives of the
other $62$ orbits there are $271$ rational character components.
Of these, $228$ have a zero simple dimension and are impossible.
For each of the other $43$, the norm of the prospective global
dimension has $17$-adic valuation zero. The original ring has a formal
codegree with minimal polynomial $t^2-272t+1088$, whose norm has
$17$-adic valuation one. Thus the signed-sphericalization
norm-divisibility condition rejects every remaining component. No
nongauge quotient can occur for an actual category.

For completeness, a gauge quotient gives sphericality of the original
category, not merely compatibility with spherical dimensions. If
$T_{ij}^k=\eta_i\eta_j\eta_kN_{ij}^k$, replace the chosen lift of $X_i$
in the sphericalization by its tensor product with the sign object
when $\eta_i=-1$. The new signed quotient has $T=N$. The sum and
difference of the multiplicities of the two lifts of $X_k$ in a tensor
product are then both $N_{ij}^k$, so the opposite lift has multiplicity
zero. These chosen lifts therefore generate a full tensor subcategory
of the sphericalization which projects tensor-equivalently onto the
original category. Transport its spherical structure along that
equivalence.

The four rational characters of the original ring have a zero entry,
so none is a spherical dimension character. The only remaining
character orbit is
\[
 (1,1,d,d,d-1,d+1),\qquad d^2-8d-1=0.
\]
One of its two real embeddings is the Frobenius--Perron character,
$d=4+\sqrt{17}$. Galois conjugation to that embedding gives a spherical
category with positive Frobenius--Perron dimensions, hence a
pseudo-unitary category. Conversely a pseudo-unitary categorification
is a complex fusion categorification.

The complete unrestricted signed enumeration has $274655$ search nodes
and is frozen in
\filename{research/new_frontier/frontier_spherical_reduction.json}.
The companion Python script recomputes all $1480$ quotients and the
complete sign orbits, verifies the original gauge orbit, checks all
$271$ nongauge character components and their arithmetic obstructions,
and reconstructs the original dimension-character orbit using exact
arithmetic. The older normalized twelve-table enumeration is retained
separately as a regression control; it is not needed for this proof.

\end{proof}


\subsection{Joint products in the former multiplicity-three frontier}
\label{sec:joint-low-product-gram}
The first general residual is still $\RR{3}{215}$. Its positive dimension
vector, computed from its multiplication tensor, is
\[
 (1,1,d,d,d-1,d+1),\qquad d=4+\sqrt{17},\qquad
 D=\FPdim(R)=136+32\sqrt{17}.
\]
A feasible induction system is not a categorification. We strengthen that
system by requiring all products of its distinguished central objects to
share one integral Gram decomposition. This retains fourth-order modular
moments as well as the restriction equations, and is more restrictive
than testing each product independently.

\begin{proposition}[Joint low-product Gram condition]
\label{prop:joint-low-product-gram}
Suppose that $\mathcal C$ is pseudo-unitary. Let
$Z_0=\mathbf1,Z_1,\ldots,Z_{s-1}$ be distinct self-dual simple objects
of $\mathcal Z(\mathcal C)$ whose restrictions $v_a=[F(Z_a)]$ are known.
Write $D=\FPdim(\mathcal C)$, $d_W=\FPdim(W)$ and
$\lambda_a(W)=S_{aW}/d_W$, using the unnormalized center matrix $S$.
For $p=(a,b)$, $1\leq a\leq b<s$, put
\[
 t_{p,e}=\frac1{D^2}\sum_W
       d_W^2\lambda_a(W)\lambda_b(W)\lambda_e(W),
 \qquad 0\leq e<s,
\]
where $\lambda_0(W)=1$. Define
\[
 Q_{(a,b),(c,d)}=
 \frac1{D^2}\sum_W d_W^2
       \lambda_a(W)\lambda_b(W)\lambda_c(W)\lambda_d(W)
       -\sum_{e=0}^{s-1}t_{(a,b),e}t_{(c,d),e}.
\]
There is a nonnegative integral matrix $B=(b_{pW})$, with columns indexed
by the remaining simple central objects, such that
\begin{equation}\label{eq:joint-low-product-gram}
 Q=BB^{\mathsf T},\qquad
 \sum_{W\notin\{Z_0,\ldots,Z_{s-1}\}}b_{pW}[F(W)]
       =v_av_b-\sum_{e=0}^{s-1}t_{p,e}v_e.
\end{equation}
In particular every $t_{p,e}$ and every entry of $Q$ is a nonnegative
integer, and $Q$ is positive semidefinite.
\end{proposition}
\begin{proof}
The Verlinde formula gives $t_{p,e}=[Z_aZ_b:Z_e]$. Self-duality makes
these rows of $S$ real and permits the displayed formulas without
conjugations. Orthogonality of the normalized matrix $S/D$ gives
\[
 [Z_aZ_b,Z_cZ_d]
   =D^{-2}\sum_W d_W^2
         \lambda_a(W)\lambda_b(W)\lambda_c(W)\lambda_d(W).
\]
Subtracting the contributions of the specified low objects leaves
$Q_{pq}=\sum_W[Z_aZ_b:W][Z_cZ_d:W]$, where the sum is over the
remaining objects. Set $b_{pW}=[Z_aZ_b:W]$ and apply the tensor functor
$F$ to each product. This proves both identities in
\eqref{eq:joint-low-product-gram}.
\end{proof}

This condition has a useful inexpensive arithmetic refinement. Group
high objects by their identical restriction vectors $v_1,\ldots,v_h$,
and let $c_j>0$ be the numbers of objects in those types. If the $v_j$
are linearly independent over $\QQ$, the restriction equation determines
uniquely the type totals $T_{pj}=\sum_{W:F(W)=v_j}b_{pW}$. The totals
must be nonnegative integers. With $T_{pj}=c_jq_{pj}+r_{pj}$,
$0\leq r_{pj}<c_j$, necessarily
\begin{equation}\label{eq:integral-variance-bound}
 Q_{pp}\ \geq\ \sum_j\bigl(c_jq_{pj}^2+
                                  r_{pj}(2q_{pj}+1)\bigr).
\end{equation}
Indeed the smallest sum of squares of $c_j$ nonnegative integers with
sum $T_{pj}$ occurs when their values differ by at most one. More
strongly, subtracting the type averages shows that
\[
 Q-T\,\operatorname{diag}(c_j^{-1})\,T^{\mathsf T}
 \quad\hbox{is positive semidefinite}.
\]
If this difference is zero, every $b_{pW}$ is forced to equal
$T_{pj}/c_j$ within its restriction type. Integrality of these averages
and both identities in \eqref{eq:joint-low-product-gram} then give a
complete exact test of this particular low-product completion.
This is a necessary condition on center data; rejecting a single
completion does not reject its input fusion ring.

For $\RR{3}{215}$ the released induction witness has six low objects
and sixteen high objects. All its low character rows are real. Their
central $S$ rows are therefore real, and
$S_{a^*,W}=\overline{S_{a,W}}$ together with invertibility of $S$ forces
these low objects to be self-dual in any realization. The automated
fourth-order test produces a $15\times15$ matrix $Q$ of rank $3$.
There are three independent high restriction types, with multiplicities
$(4,8,4)$, and the type-variance matrix is exactly zero. Every forced
average is a nonnegative integer. Expanding the three types into all
sixteen formal objects gives an explicit integral matrix $B$ satisfying
both equations exactly.

The calculation and its frozen certificate are
\filename{research/m8_frontier/low_product_gram.py} and
\filename{research/m8_frontier/low_product_gram_certificate.json}.
The program reconstructs the tensor, character table, inherited
230-equation witness, every cubic and fourth-order moment, all product
restrictions, and the expanded Gram identities. Thus the stronger test
survives and contributes zero general exclusions. On its own it supplies
no complete center fusion tensor or twists. The following subsection
extends the witness to those additional data. A further source-specific
argument then obstructs lifting its displayed map. These older necessary
systems alone do not settle general existence; the stabilizer-rotation
criterion in Theorem~\ref{thm:frontier-noncategorifiable} supplies the
general exclusion.

\subsection{An exact center-ring completion of the multiplicity-three data}
\label{sec:cyclic-eight-center}
The surviving low-product witness extends substantially further than the
Gram conditions alone: it admits a full modular datum from a known unitary
category, together with a positive based-ring map to $\RR{3}{215}$.
This is still weaker than a tensor functor or a categorification of the
rank-six ring.

Let $\mathcal D$ be the $\ZZ/2\ZZ$-de-equivariantization of the generalized
Haagerup category for $\ZZ/8\ZZ$ in
\cite[Theorem~6.8 and Example~6.10]{GIcenters}. Put $d=4+\sqrt{17}$ and
$D=136+32\sqrt{17}$. We use the source's labels
$\mathbf1,\pi,\varphi_\pm,\sigma^g_k,\psi_{\epsilon_1,\epsilon_2},\xi_j$
for the $22$ simples of $Z(\mathcal D)$.
Here the six $\sigma$ labels are $(g,k)=(0,1),(2,1),(1,0),\ldots,(1,3)$,
the four $\psi$ labels have $\epsilon_1,\epsilon_2\in\{\pm1\}$, and
$1\le j\le8$.

\begin{proposition}[Exact compatibility with a known center]
\label{prop:cyclic-eight-center}
There is a unital, duality-preserving homomorphism
\[
 f:K_0(Z(\mathcal D))\longrightarrow \RR{3}{215}
\]
which sends every simple basis element to a nonnegative integral combination
of the target basis. Its values are the vectors in the following table.
The target basis has dimensions $(1,1,d,d,d-1,d+1)$.
\begin{center}
\small
\begin{tabular}{@{}llll@{}}
\toprule
Central label & Dimension & Twist & $f$-vector\\
\midrule
$\mathbf1$ & $1$ & $1$ & $(1,0,0,0,0,0)$\\
$\pi$ & $8d+1$ & $1$ & $(1,0,2,2,2,2)$\\
$\sigma^0_1$ & $8d+2$ & $1$ & $(1,1,1,3,2,2)$\\
$\sigma^2_1$ & $8d+2$ & $1$ & $(1,1,3,1,2,2)$\\
$\varphi_+$ & $4d+1$ & $1$ & $(1,0,0,0,2,2)$\\
$\varphi_-$ & $4d+1$ & $1$ & $(1,0,2,2,0,0)$\\
$\sigma^1_k$, $0\le k<4$ & $8d+2$ & $\zeta_8^{1+2k}$ & $(0,0,2,2,1,3)$\\
$\xi_j$, $1\le j\le8$ & $8d$ & $\zeta_{17}^{3j^2}$ & $(0,0,2,2,2,2)$\\
$\psi_{\epsilon_1,\epsilon_2}$ & $4d+1$ & $-1$ & $(0,1,1,1,1,1)$\\
\bottomrule
\end{tabular}
\end{center}
Every simple object of this center has second Frobenius--Schur indicator
$+1$.
\end{proposition}
\begin{proof}
We specialize the cited modular-data formulas to $m=2$ and reconstruct the
complete center fusion tensor by the Verlinde formula. The exact certificate
stores its $22^3$ structure constants, the normalized matrix $s$, its twists,
and the displayed $22\times6$ matrix $F$.
All entries of $136s$ belong to
$\ZZ[x]/(\Phi_{136}(x))$, where
\[
 \Phi_{136}(x)=\sum_{j=0}^{16}(-x^4)^j,
 \qquad x=\zeta_{136}.
\]
The checker verifies symmetry, $s^2=I$ and
$sts=t^{-1}st^{-1}$ exactly, where $t$ is the diagonal twist matrix.
For every $a,b,\ell$ it verifies
\[
 \sum_c N_{ab}^{c}s_{c\ell}s_{0\ell}
       =s_{a\ell}s_{b\ell}.
\]
Together with invertibility of $s$, these $10\,648$ identities certify the
Verlinde tensor without numerical rounding. Its coefficients are
nonnegative integers; the unit, duality and associativity identities are
checked independently. For every $a,b,i$ the checker then verifies
\[
 \sum_c N_{ab}^{c}F_{ci}
 =\sum_{j,k}F_{aj}F_{bk}N_{jk}^{i}(\RR{3}{215}).
\]
These $2\,904$ identities prove the asserted ring homomorphism. Finally it
checks the second-indicator formula
\[
 \nu_2(a)=\sum_{i,j}N_{ij}^{a}s_{0i}s_{0j}
                     (t_i/t_j)^2=1
\]
in the same cyclotomic field.
\end{proof}

The proof is reproduced by
\filename{research/center_continuation/exact_cyclic8_center.py}; it remains
valid under Python's optimization flag. The frozen data are
\filename{cyclic8_center_tensor.json} and
\filename{cyclic8_modular_data.json} in the same directory.
Thus the frontier is compatible with a complete known central fusion ring,
individual twists, modular relations and all second indicators.
These compatibility checks do not produce an algebra or a central tensor
functor. In fact, Proposition~\ref{prop:cyclic-eight-lift-obstruction}
proves that this specified map cannot be the central forgetful functor
under an equivalence with the actual known center used here. This does
not identify every modular category with the same datum. The general
nonexistence result instead follows from the stabilizer-rotation
obstruction in Theorem~\ref{thm:stabilizer-rotation-frontier}, which
applies to every possible center. The numerical compatibility above
therefore illustrates why a fusion tensor and modular datum alone do
not provide a categorification.


\subsection{The displayed restriction map does not lift categorically}
\label{sec:cyclic-eight-lift-obstruction}
The distinction between a positive ring map and a central tensor functor
is essential in Proposition~\ref{prop:cyclic-eight-center}.

\begin{proposition}[Obstruction for the specified cyclic-eight data]
\label{prop:cyclic-eight-lift-obstruction}
Let $\mathcal D$ be the actual cyclic-eight quotient category used in
Proposition~\ref{prop:cyclic-eight-center}, and let $f$ be its displayed
positive based-ring map. There is no fusion category $\mathcal C$ with
based ring $\RR{3}{215}$ and a braided equivalence
$Z(\mathcal D)\simeq Z(\mathcal C)$ under which the central forgetful
functor induces $f$. The same holds after dimension- and twist-preserving
permutations of the displayed central labels.
\end{proposition}
\begin{proof}
Write $d=4+\sqrt{17}$, so $d^2=8d+1$. The simple objects of
$\mathcal D$ are $\alpha^j,\alpha^j\rho$ for $j\in\ZZ/4\ZZ$, with
\[
 \rho\alpha=\alpha^{-1}\rho,
 \qquad \rho^2=1+2\sum_j\alpha^j\rho.
\]
The category has connected separable algebras $1+\alpha^j\rho$ for
every $j$; these come from the generalized Haagerup construction and
descend under the quotient tensor functor
\cite[Section~2.4 and Example~4.4]{GIcenters}.

Suppose the asserted equivalence and central functor exist, and let
$\mathcal M$ be the corresponding $\mathcal D$--$\mathcal C$ Morita
bimodule. The actual source restriction matrix $F_{\mathcal D}$ is
recorded in
\filename{research/center_continuation/cyclic8_canonical_restriction.json}.
Its complete tensor compatibility is verified against the same center
tensor as $f$. The trace identity given by the cross Gram matrix
$F_{\mathcal D}^{t}F$ yields
\[
 |\Irr(\mathcal M)|=6,\quad
 \operatorname{Tr}(L_\alpha)=0,\quad
 \operatorname{Tr}(L_{\alpha^2})=2,\quad
 \operatorname{Tr}(L_{\alpha^j\rho})=8,\quad
 \operatorname{Tr}(R_g)=2.
\]
Here $g$ is the nontrivial invertible of $\mathcal C$. The two columns
of $F$ used in these identities are constant on each class of central
simples having equal dimension and twist, so the permitted relabelings
do not change the argument.

The permutation $L_\alpha$ has one four-cycle and one two-cycle.
Symmetry and $L_\rho L_\alpha=L_\alpha^{-1}L_\rho$ put $L_\rho$ in
a reverse-circulant four-by-four block, an alternating cross block, and
a lower block $\left(\begin{smallmatrix}u&v\\v&u\end{smallmatrix}\right)$.
The trace equations give
$u=4-r_0-r_2$, $v=4-r_1-r_3$ and
$\operatorname{Tr}(L_\rho^2)=70$.
The exact enumeration in
\filename{research/new_frontier/cyclic8_rank6_modules.py}
checks all $2025$ bounded integer assignments. Exactly eight remain,
four of each type in the table below. Both cross entries equal one.
The squared dimensions use Morita normalization, so their sum is
$\operatorname{FPdim}(\mathcal D)=8+32d$.
\begin{center}
\begin{tabular}{@{}ccc@{}}
\toprule
Upper row sum & Four-cycle objects & Two-cycle objects\\
\midrule
$1$ & $1+d$ & $2+14d$\\
$7$ & $1+7d$ & $2+2d$\\
\bottomrule
\end{tabular}
\end{center}
The unique involutory permutation of trace two commuting with
$L_\alpha$ shifts its four-cycle by two and fixes its two-cycle.
This determines $R_g$.

The algebra $A=1+g$ exists because $g$ fixes a simple object. Indeed that
object gives a one-object module over the pointed subcategory generated
by $g$, trivializing its associator. The exact full Morita calculation
\filename{research/new_frontier/verify_frontier_full_morita.py}
shows that the based ring of
$\mathcal C'=A\text{-bimod}(\mathcal C)$ is again $\RR{3}{215}$.
It does not identify $\mathcal C$ and $\mathcal C'$ categorically.
In the second row of the table, pass to right $A$-modules in
$\mathcal M$. Each fixed simple has two module structures, whose
squared Morita dimensions are $(2+2d)/2=1+d$; each free pair gives one
simple of squared dimension $4(1+7d)/2=2+14d$.
Thus in either case there is a Morita bimodule, with right based ring
$\RR{3}{215}$, containing a simple $\iota$ with
$\iota\iota^*=1+\alpha^h\rho$. Rename $\alpha^h\rho$ as $\rho$.

For clarity, this reduction does not assume that every algebra
$1+\rho$ has the same module action. Its four $\alpha$-translates are
distinct. The upper $\rho$ block and its diagonal square equation force
both cross entries to be one. The remaining lower block has
$u+v=7$ and $|u-v|=1$. If $u=4$, the object $\iota^*\kappa$ in the
right category would have dimension $4d$ and endomorphism dimension
$5$. It has no invertible summands since $\iota$ and $\kappa$ have
different dimensions. Writing its other multiplicities as $a,b,c,e$
gives $e=c$, $a+b+2c=4$ and squared norm $a^2+b^2+2c^2$.
The possible norms are $4,6,8,10,16$, never $5$. Hence $u=3$, $v=4$.

Choose a Morita object $\lambda$ in a second module category over
$\mathcal D$ with $\lambda\lambda^*=1+\alpha\rho$.
In the composed bimodule, Frobenius reciprocity gives two distinct
simples $\beta_0=\lambda^*\iota$ and
$\beta_2=\lambda^*\alpha^2\iota$, each of dimension $d+1$.
The object $\Gamma=\lambda^*\kappa$ has dimension $4d$, endomorphism
dimension $5$, and contains both $\beta$'s once. Its three other simple
summands have total dimension $2(d-1)$. One of them, $\gamma$, therefore
has dimension at most $2(d-1)/3$.

All dimensions in this composed bimodule lie in $K=\QQ(\sqrt{17})$:
the simple Perron eigenspace of the sum of the right action matrices
is defined over $K$, and the known dimension $d+1$ fixes its scale.
Thus the connected internal end $B=\gamma^*\gamma$ has dimension a
square in $K$, at most $L=4(d-1)^2/9$. It cannot be the unit algebra;
otherwise the bimodule would be the regular module and its two distinct
simples of dimension $d+1$ would contradict the dimensions of
$\RR{3}{215}$.

The unit has multiplicity one in $B$ and the other invertible has
multiplicity at most one. The complete list of dimensions at most $L$ is
\[
 1,2,d,d+1,d+2,d+3,
 2d-1,2d,2d+1,2d+2,2d+3,2d+4,3d-2.
\]
Here $2$ is not a square in $K$, and the field norms of the other
nontrivial entries are respectively
\[
 -1,8,19,32,-19,-4,13,32,53,76,-53,
\]
none a nonnegative rational square. This contradicts the existence of
$B$. The exhaustive bound and these norms are verified by
\filename{research/new_frontier/cyclic8_small_algebra_dimensions.py}.
This extends the second-algebra argument of
\cite[Lemmas~4.15--4.16]{GISAH}; the additional possibility $3d-2$,
omitted from that paper's small-dimension list, has norm $-53$ and
does not affect the contradiction.
\end{proof}

The source category and specified induction data in this proposition
are substantive hypotheses. An abstract modular datum or an integral
Verlinde tensor does not identify a braided category with $Z(\mathcal D)$.
Consequently this proposition alone does not resolve general
categorifiability. The independent stabilizer-rotation argument in
Theorem~\ref{thm:stabilizer-rotation-frontier} supplies that general
exclusion without identifying an unknown center with $Z(\mathcal D)$.

